\documentclass[11pt]{exam}
\usepackage[empty]{fullpage}
\usepackage{amsmath,amssymb}
\usepackage{colortbl}
\usepackage{subfig}
\usepackage{multicol}
\usepackage{tikz,pgfplots}
\usepgfplotslibrary{statistics}
\usepackage{color,comment,soul}
\usepackage{pdfpages}

\newcommand{\blank}[1]{\underline{\hspace*{#1}}}
\newcommand{\ds}{\displaystyle}
\newcommand{\on}{\operatorname}
\newcommand{\Var}{\operatorname{Var}}

\addpoints
\vqword{Problem}
\vpword{Maximum Points}
\vsword{Your Score}
\gradetablestretch{1.8}
%\printanswers

\begin{document}
\graphicspath{{../../Semesters12-19/Spring18/Math222/}}

\section*{Inference Formulas \hfill Math 222}
\begin{description}


%\item[Probability Rules]
%$$P(A \text{ or } B) = P(A) + P(B) - P(A \text{ and } B)$$
%$$P(A \text{ and } B) = P(A) P(B|A)$$
%\vfill
%
%
%
%\item[Expected Value \& Variance] For a discrete random variable $X$ with outcomes $x_k$ and corresponding probabilities $p_k$.  
%$$E(X) = \mu = \sum p_k x_k$$
%$$\Var(X) = \sigma^2 = \sum p_k (x_k - \mu)^2$$
%\vfill
%
%\item[Expected Value \& Variance Rules] For any random variables $X$ and $Y$.
%
%\begin{center}
%\begin{tabular}{ccc}
%$E(X+Y) = E(X)+E(Y)$ & &  $E(cX) = cE(X)$ \\
%$\Var(X+Y) = \Var(X)+\Var(Y)^*$ & \hspace*{0.5in} & $\Var(cX) = c^2 \Var(X)$  
%\end{tabular}
%\end{center}
%$^*$\textit{Variance is only additive if $X$ and $Y$ are independent.}
%\vfill
%
%
%%\item[Mean \& Variance of a Sample Mean]
%%$$E(\bar{x}) = \mu_{\bar{x}} = \mu$$
%%$$\Var(\bar{x}) = \sigma^2_{\bar{x}} = \frac{\sigma^2}{{n}}$$
%%\vfill
%
%\item[Mean \& Variance for the Binomial Distribution] If $X \sim \operatorname{Binom}(n,p)$, then
%$$E(X) = \mu = np$$
%$$\Var(X) = \sigma^2 = p(1-p)n$$
%\vfill
%
%%\item[One Sample Inference for Means] ~
%%$$\bar{x} \pm t^* \frac{s}{\sqrt{n}} \hspace*{0.5in} \ds t = \frac{\bar{x} - \mu}{s/\sqrt{n}}$$
%%\begin{center}
%%\begin{tabular}{|c|c|}
%%\hline
%%$\ds \bar{x} \pm t^* \frac{s}{\sqrt{n}}$ & $\ds t = \frac{\bar{x} - \mu}{s/\sqrt{n}}$ \\ \hline
%%$\ds (\bar{x}_1 - \bar{x}_2) \pm t^* \sqrt{ \frac{s_1^2}{n_1} + \frac{s_2^2}{n_2} }$ & $\ds t = \frac{\bar{x}_1 - \bar{x}_2}{ \sqrt{ \frac{s_1^2}{n_1} + \frac{s_2^2}{n_2}} }$ \\ \hline
%%\end{tabular}
%%\end{center}
%
%
%%
%%\item[Two Sample Inference for Means]
%%\[
%%(\bar{x}_1 - \bar{x}_2) \pm t^* \sqrt{ \frac{s_1^2}{n_1} + \frac{s_2^2}{n_2} } \hspace*{0.5in} t = \frac{\bar{x}_1 - \bar{x}_2}{ \sqrt{ \frac{s_1^2}{n_1} + \frac{s_2^2}{n_2}} }
%%\]
%%
%%\item[One Sample Inference for Proportions]
%%$$\hat{p} \pm z^* \sqrt{\frac{\hat{p}(1-\hat{p})}{n}} \hspace*{0.25in} z = \frac{\hat{p}-p_0}{\sqrt{\frac{p_0(1-p_0)}{n}}}$$
%%
%%
%%\item[Two Sample Inference for Proportions]
%%\[
%%(\hat{p}_1 - \hat{p}_2)  \pm z^* \sqrt{ \frac{\hat{p}_1(1-\hat{p}_1)}{n_1}+\frac{\hat{p}_2(1-\hat{p}_2)}{n_2} } \hspace{0.25in} z = \frac{\hat{p}_1 - \hat{p}_2}{\sqrt{\hat{p}(1-\hat{p})\left( \frac{1}{n_1} + \frac{1}{n_2} \right)}}
%%\]


\item[General form of a confidence interval] $\text{estimate} \pm \text{margin of error}$, where
$$\text{margin of error} = (\text{critical value}) \times SE_{\text{estimate}}$$
\vfill

\item[General form of a test statistic]
$$ \frac{\text{statistic} - \text{hypothesized valued}}{SE_{\text{statistic}}}$$
\vfill

\item[Standard errors$^*$] ~
\begin{center}
\begin{tabular}{lcl} 
$\ds SE_{\bar{x}} = \frac{s}{\sqrt{n}}$ & ~~~~ & $\ds SE_{\hat{p}} = \sqrt{\frac{\hat{p}(1-\hat{p})}{n}}$ \\
$\ds SE_{\bar{x}_1 - \bar{x}_2} = \sqrt{\frac{s_1^2}{n_1}+\frac{s_2^2}{n_2}}$ &  & $\ds SE_{\hat{p}_1-\hat{p}_2} = \sqrt{\frac{\hat{p}_1(1-\hat{p}_1)}{n_1}+\frac{\hat{p}_2(1-\hat{p}_2)}{n_2}}$ \\
\end{tabular}
\end{center}
$^*$\textit{When testing a null hypothesis about proportions, it is better to use the hypothesized population proportion $p_0$ rather than $\hat{p}$ in the formula for standard error for one-sample tests, and it is better to used the pooled proportion $\hat{p}$ rather than either individual sample proportion $\hat{p}_1$ or $\hat{p}_2$ in the formula for standard error in two-sample tests. }
\vfill


\item[Chi-squared formulas] ~
$$\ds \chi^2 = \sum \frac{(\textit{observed count}-\textit{expected count})^2}{\textit{expected count}}$$ 

\item[Expected counts for 2-way tables] ~ 
$$E_{ij} = \frac{\textit{row total} \times \textit{column total}}{\textit{table total}}$$ 

\item[Least squares regression line] ~
$$\hat{y} = b_0+b_1 x, ~\text{  where  }~ b_1 = r \frac{s_y}{s_x} ~\text{  and  }~ b_0 = \bar{y} - b_1 \bar{x}$$

\item[Inference about regression] 
$$SE_{b_1} = \frac{s}{s_x \sqrt{n - 1}} = \frac{r \sqrt{n-2}}{1-r^2}$$ 
$$SE_{\mu_y} = s \sqrt{\frac{1}{n} + \frac{(x^* - \bar{x})^2}{s_x^2 (n-1)}} = \sqrt{\frac{s^2}{n} + (x^* - \bar{x})^2 SE_{b_1}^2}.$$
$$SE_{\hat{y}} = s \sqrt{1 + \frac{1}{n} + \frac{(x^* - \bar{x})^2}{s_x^2 (n-1)}} = \sqrt{s^2 + \frac{s^2}{n} + (x^* - \bar{x})^2 SE_{b_1}^2}.$$
%$$SE_{\hat{y}} = s \sqrt{1 +  \frac{1}{N} + \frac{ b_1^2 (x-\bar{x})^2}{MSM}}$$ 
%$$\ds t = \frac{b_1}{SE_{b_1}} = \frac{r\sqrt{n-2}}{\sqrt{1-r^2}} = \sqrt{F} = \sqrt{\frac{MSM}{MSE}}$$
where $s = \sqrt{MSE}$ is the standard error of the residuals.
Recall that $r^2 = \frac{SSM}{SST}$ and $MST = s_y^2$. With that you can fill in the whole ANOVA table.\\


\item[Inference in ANOVA] Use the pooled standard deviation $s_p = \sqrt{MSE}$ as the best guess for $\sigma$.   
Recall also that $SSG = \sum_{i = 1}^I N_i (\bar{x}_i - \bar{x})^2$ and $SSE = \sum_{i = 1}^{I} (N_i-1) s_i^2$. \\
%For a contrast $C=\sum_i a_i \bar{x}_i$ the standard error is $SE_C = s_p \sqrt{\sum \frac{a_i^2}{N_i} }$.

%\begin{center}
%\begin{tabular}{lcl}
%Testing statistic  $H_0: \beta_1 = 0$ & & Residual standard error \\[0.2in] 
%$\ds t = \frac{b_1}{SE_{b_1}} = \frac{r\sqrt{n-2}}{\sqrt{1-r^2}}$ & ~~~~ & $\ds s = \sqrt{\frac{\sum(y_i - \hat{y}_i)^2}{n-2}}$ \\[0.2in]
%$\ds SE(b_0) = s \sqrt{\frac{1}{n}+\frac{\bar{x}^2}{\sum (x_i-\bar{x})^2}}$ & ~~~~ & $\ds SE(b_1) = \frac{s}{\sqrt{\sum(x_i-\bar{x})^2}}$ \\[0.2in]
%$\ds SE(\hat{\mu}_y) = \sqrt{SE(b_0)^2 + (x^*)^2 SE(b_1)^2}$ &  & $\ds SE(\hat{y}) = \sqrt{SE(\hat{\mu}_y)^2+s^2} $ \\
%\end{tabular}
%\end{center}

\end{description}

\end{document}
